% Notes of the Montreal seminar, 1970 (cote 129, typed part): archive cote 129, site page 11 (PDF page 12).
% typescript. Draft compared: none.
% AI-transcribed by Claude Opus 5.5 (claude-opus-5-5, adaptive thinking, effort high), Message Batches API, batch mtl129_full (msgbatch_01AGnDKHwy6ATCM7ehUa6263), 2026-11-13; not yet checked by eye.
\scanpage{129/11}
\diagram{\begin{tikzcd}
X \arrow[r, "\mathrm{ff}_X"] \arrow[d, "g"'] & X \arrow[d, "g"'] \\
Y \arrow[r, "\mathrm{ff}_Y"] & Y
\end{tikzcd}}

On obtient donc un endomorphisme $\mathrm{ff}$ du foncteur identique de $\mathrm{Sch}/\mathbb{F}_p$.

5.2. Soit $S$ un schéma fixé et $X$ un S-schéma. On note $X^{(p/S)}$
ou simplement $X^{(p)}$ le S-schéma image réciproque de $X$ par le changement
de base $\mathrm{ff}_S : S \to S$.

Le diagramme ci-dessus étant commutatif, il existe une flèche
unique $\mathrm{ff}_{X/S}$ telle que le diagramme suivant soit commutatif:

\diagram{\begin{tikzcd}
X \arrow[dr, "\mathrm{ff}_{X/S}"'] \arrow[drr, "\mathrm{ff}_X"] \arrow[ddr] & & \\
& X^{(p)} \arrow[r] \arrow[d] & X \arrow[d] \\
& S \arrow[r, "\mathrm{ff}_S"] & S
\end{tikzcd}}

On appellera le morphisme $\mathrm{ff}_{X/S}$ ainsi défini le \uline{morphisme de Frobenius}
\uline{relatif de} $X$ \uline{au-dessus de} $S$.

L'application $X \mapsto \mathrm{ff}_{X/S}$ est un morphisme fonctoriel du foncteur
identique de $\mathrm{Sch}/S$ dans le foncteur $X \mapsto X^{(p/S)}$. Ce dernier foncteur
commute aux limites projectives finies, donc si $X$ est un S-groupe,
$X^{(p/S)}$ l'est aussi et, $\mathrm{ff}_{X/S}$ étant fonctoriel, est un morphisme de groupes.

La construction de $X^{(p/S)}$ commute aux changements de base,
c'est-à-dire que pour tout morphisme $T \to S$, on a

\mspage{\hand{11}}
